% ============================================================================
%  sections/part3a_utility_core.tex  --  Part III, Section 7
%  "The acquisition objective" (utility core).
%  \input-ready fragment: no preamble/documentclass. Uses the shared macros
%  from preamble.tex. Source: notes/06_utility_core.md; ground truth = the
%  current code (acquisition.py, utils.py). The deep proofs (conditioning,
%  divergence, predictive-mean equivalence) live in the sibling Part III-deep
%  section and are only cited here.
% ============================================================================

\section{The acquisition objective}\label{sec:utility}

Active learning closes the loop: after each spike-count observation the model is
refit, and the next natural image to display is the one whose observation is
expected to teach us the most about the neuron's tuning. ``Teach us the most'' is
made precise as an \emph{information gain} --- an expected reduction in the entropy
of what we do not yet know about the retinal ganglion cell's firing behaviour.
This section builds that utility on the generative model and variational posterior
of Parts~I--II: the Poisson-lognormal predictive of the spike count, its entropy,
and the two utilities the codebase actually computes.

Throughout, the candidate (query) image is $x^{*}$; a spike count is the response
random variable $R$ with realizations $r$ (this is the spike count written $\spk$
elsewhere in the document --- we keep it as $r$ here to match the count-domain
code variables $\rmax$ and the Laplace mode $\gmode$). The neuron's mean firing
rate to image $x$ is $\rate(x)=\exp(\gain\,\lat(x)+\latz)$ with latent
$\lat(x)$, gain $\gain$, and offset $\latz$; the spike count is
$R\mid\rate(x)\sim\Poi(\rate(x))$. It is convenient to work with the
\emph{log-firing rate} $\lgr(x)=\gain\,\lat(x)+\latz$, so $\rate=e^{\lgr}$, whose
GP posterior at any image is Gaussian with the log-rate moments
%
\begin{equation}
  \mug \;=\; \gain\,\pmean(x^{*})+\latz,
  \qquad
  \varg \;=\; \gain^{2}\,\pvar(x^{*}),
  \label{eq:logratemoments}
\end{equation}
%
the affine push-forward of the latent posterior moments $\pmean(x^{*})$,
$\pvar(x^{*})$ read off the variational posterior~\eqref{eq:predmoments}.

\begin{remark}[Two utilities, one idea --- and which one ships]\label{rem:two-utilities}
The codebase carries \emph{two} acquisition functions. Both express the same idea
--- \emph{utility$(x^{*})=$ information the observation $R$ at $x^{*}$ gives about
the neuron's tuning $=$ a drop in entropy} --- and they differ only in
\emph{which} uncertainty they score (Table~\ref{tab:two-utilities}). The
\textbf{standard utility} $\Ustd$ scores the firing rate \emph{at $x^{*}$ itself};
the \textbf{distribution-aware utility} $\Uda$ scores the firing rates at the
natural images $x\sim\px$ we ultimately want to predict. The production
closed-loop experiment selects images with $\Ustd$ (verified:
\texttt{run\_active\_loop.py:241} calls \texttt{standard\_utility}). $\Uda$ is the
theoretically richer object derived in the internal write-ups, but in the current
code it is a research instrument used only under \texttt{investigations/} --- and
it carries a norm-exploitation divergence (\S\ref{subsec:landscape},
Thm.~\ref{thm:divergence}) that the finite image pool masks in production. We label
$\Ustd$ = production, $\Uda$ = research, throughout.
\end{remark}

\begin{table}[t]
  \centering
  \small
  \begin{tabular}{@{}lll@{}}
    \toprule
      & \textbf{Standard} $\Ustd$ (production) & \textbf{Distribution-aware} $\Uda$ (research)\\
    \midrule
    Formula      & $\Hmarg(x^{*})-\E[\Hnoise(x^{*})]$
                 & $\Hmarg(x^{*})-\E_{x\sim\px}[\Hcond(x^{*})]$\\
    As MI        & $I\!\big(\rate(x^{*});R\mid x^{*},\mathcal D\big)$
                 & $I\!\big(\rate(X);R\mid X,x^{*},\mathcal D\big),\ X\sim\px$\\
    Scores       & rate \emph{at $x^{*}$ itself}
                 & rates at \emph{natural images} $x\sim\px$\\
    Needs        & marginal moments at $x^{*}$
                 & marginal moments \emph{and} cross-cov.\ to $\px$\\
    Code         & \texttt{standard\_utility}$\to$\texttt{nd\_utility\_new}
                 & \texttt{distribution\_aware\_utility}\\
    Where        & \texttt{run\_active\_loop.py:241} (real experiment)
                 & \texttt{investigations/} only\\
    \bottomrule
  \end{tabular}
  \caption{The two utilities. Both are ``information gain $=$ entropy reduction'',
  about different targets. $\Ustd$ ships; $\Uda$ is a research variant.}
  \label{tab:two-utilities}
\end{table}

% ---------------------------------------------------------------------------
\subsection{Utility as expected entropy reduction and mutual information}
\label{subsec:mi}

Fix a candidate $x^{*}$ and the data $\mathcal D=\{(x_i,r_i)\}$. The utility is the
mutual information between the response $R$ we would record at $x^{*}$ and the
firing behaviour we want to learn. What ``firing behaviour'' means is exactly the
difference between the two utilities.

\paragraph{The distribution-aware objective.}
Let $Z:=(X,\rate(X))$ with $X\sim\px$ a natural image and $\rate(X)$ its mean
firing rate. Choosing $x^{*}$ to maximize the information $R$ carries about $Z$,
%
\begin{equation}
  \Uda(x^{*}) \;:=\; I\!\big(Z;R\mid x^{*},\mathcal D\big)
             \;=\; I\!\big((X,\rate(X));\,R\mid x^{*},\mathcal D\big).
  \label{eq:da-objective}
\end{equation}
%
The chain rule for mutual information splits this into a term for \emph{which}
natural image we imagine and a term for its \emph{rate}:
%
\begin{equation}
  I\!\big((X,\rate(X));R\mid x^{*},\mathcal D\big)
  \;=\;
  \underbrace{I(X;R\mid x^{*},\mathcal D)}_{=\,0}
  \;+\;
  I\!\big(\rate(X);R\mid X,x^{*},\mathcal D\big).
  \label{eq:mi-chain}
\end{equation}
%
The first term vanishes: \emph{which} natural image $X$ we intend to ask about is
independent of the response $R$ recorded at $x^{*}$, given $x^{*}$ and
$\mathcal D$. Hence
%
\begin{equation}
  \Uda(x^{*})
  \;=\; I\!\big(\rate(X);R\mid X,x^{*},\mathcal D\big)
  \;=\; \E_{x\sim\px}\!\big[\,I\!\big(\rate(x);R\mid x,x^{*},\mathcal D\big)\big].
  \label{eq:da-expectation}
\end{equation}

\paragraph{Entropy decomposition.}
For a fixed natural image $x$, the per-image mutual information decomposes into a
marginal-entropy term minus an expected conditional-entropy term,
%
\begin{equation}
  I\!\big(\rate(x);R\mid x,x^{*},\mathcal D\big)
  \;=\;
  H(R\mid x^{*},\mathcal D)
  \;-\;
  \E_{\rate(x)}\!\big[\,H(R\mid x^{*},\rate(x),\mathcal D)\big],
  \label{eq:mi-entropy-split}
\end{equation}
%
using (i) $H(R\mid x,x^{*},\mathcal D)=H(R\mid x^{*},\mathcal D)$ --- the marginal
of $R$ at $x^{*}$ does not depend on which natural image we intend to predict ---
and (ii) that conditioning on the latent value $\rate(x)$ leaves the expected
entropy $\E_{\rate(x)}[H(R\mid x^{*},\rate(x),\mathcal D)]$. Averaging
\eqref{eq:mi-entropy-split} over $x\sim\px$ and naming the two pieces gives the
distribution-aware utility as a difference of entropies,
%
\begin{equation}
  \Uda(x^{*}) \;=\; \Hmarg(x^{*}) \;-\; \E_{x\sim\px}\!\big[\Hcond(x^{*})\big],
  \label{eq:Uda}
\end{equation}
%
with
%
\begin{align}
  \Hmarg(x^{*}) &:= H(R\mid x^{*},\mathcal D), \label{eq:def-Hmarg-cond}\\
  \Hcond(x^{*}) &:= \E_{x\sim\px,\ \lat(x)\sim q}\!\big[\,H(R\mid x^{*},\lat(x),\mathcal D)\big].
  \label{eq:def-Hcond}
\end{align}
%
By the symmetry of mutual information, $\Uda(x^{*})$ equals the reduction in
uncertainty about the firing rates $\{\rate(x):x\sim\px\}$ produced by observing
$R$ at $x^{*}$ --- the precise sense of ``how much does showing $x^{*}$ teach us
about the cell's tuning on natural images.''

\begin{remark}[Two expectations inside $\Hcond$]\label{rem:double-exp}
$\Hcond$ in~\eqref{eq:def-Hcond} carries \emph{two} expectations: over natural
images $x\sim\px$ and over the posterior latent $\lat(x)\sim q$. The latent is
integrated because $\rate(x)=\exp(\gain\lat(x)+\latz)$ is itself uncertain --- only
the posterior of $\lat(x)$ is known, not its value. Replacing the inner
expectation by an evaluation at the mean $\lat(x)=\pmean(x)$ is \emph{biased}; the
necessity of the inner sampling is the subject of \S\ref{subsec:why-sample}.
\end{remark}

The \emph{standard} utility instead scores the rate at the query point itself. It
is the special case in which the ``natural image we predict'' is $x^{*}$, so
$\Uda$'s outer expectation collapses and $\Hcond$ becomes the aleatoric Poisson
noise entropy at $x^{*}$; we develop it in \S\ref{subsec:standard} once $\Hmarg$
is in hand.

% ---------------------------------------------------------------------------
\subsection{The marginal spike-count entropy \texorpdfstring{$\Hmarg$}{H\_marg}}
\label{subsec:hmarg}

Both utilities share the marginal predictive entropy of the spike count at
$x^{*}$,
%
\begin{equation}
  \Hmarg(x^{*})
  \;=\; -\sum_{r=0}^{\infty} p(r\mid x^{*},\mathcal D)\,\log p(r\mid x^{*},\mathcal D),
  \label{eq:Hmarg}
\end{equation}
%
where the marginal predictive integrates the Poisson likelihood against the
Gaussian posterior of the log-rate $\lgr(x^{*})\sim\Gauss(\mug,\varg)$
(cf.~\eqref{eq:logratemoments}):
%
\begin{equation}
  p(r\mid x^{*},\mathcal D)
  \;=\; \int \Poi(r\mid e^{\lgr})\,\Gauss(\lgr\mid\mug,\varg)\,d\lgr
  \;=\; \frac{1}{r!}\int \exp\!\big(r\lgr - e^{\lgr}\big)\,\Gauss(\lgr\mid\mug,\varg)\,d\lgr .
  \label{eq:marg-predictive}
\end{equation}
%
This integral has no closed form. The code evaluates it by a \textbf{Laplace
expansion} of the integrand about its mode $\gmode$ (the mode of the
log-integrand $r\lgr - e^{\lgr} - (\lgr-\mug)^{2}/2\varg$):
%
\begin{equation}
  \log p(r\mid x^{*},\mathcal D)
  \;\approx\;
  \gmode\,r - e^{\gmode}
  - \frac{(\gmode-\mug)^{2}}{2\varg}
  - \tfrac12\log\!\big(1+\varg\,e^{\gmode}\big)
  - \log r! .
  \label{eq:laplace-logp}
\end{equation}

\begin{lemma}[Closed-form Laplace mode via Lambert~$W$]\label{lem:mode}
The stationarity condition $r-e^{\lgr}-(\lgr-\mug)/\varg=0$ rearranges to
$\lgr=r\varg+\mug-\varg e^{\lgr}$. Substituting $w=\varg e^{\lgr}$ gives
$w+\log w=\log\varg+r\varg+\mug$, i.e.\ $w=\LW\!\big(\varg\,e^{\,r\varg+\mug}\big)$
with $\LW$ the principal branch of the Lambert~$W$ function. Hence the mode is
closed-form,
%
\begin{equation}
  \gmode \;=\; r\varg+\mug-\LW\!\big(\varg\,e^{\,r\varg+\mug}\big).
  \label{eq:laplace-mode}
\end{equation}
\end{lemma}

The infinite sum in~\eqref{eq:Hmarg} decays fast (retinal firing rates are low)
and is truncated at $\rmax$ (\S\ref{subsec:landscape} handles the truncation
artifact). At very small variance the Laplace form is unstable, and the code falls
back to the exact Poisson log-pmf $\log p(r)=r\mug-e^{\mug}-\log r!$ via a
differentiable branch (threshold $\varg<10^{-6}$). The mode is computed in
log-space --- $\LW(e^{y})$ with $y=\log\varg+r\varg+\mug$, Newton-iterated with a
custom backward $dW/dy=W/(1+W)$ --- so no $e^{r\varg+\mug}$ is ever formed
(overflow-safe, gradient-safe).

% ---------------------------------------------------------------------------
\subsection{The standard (production) utility}
\label{subsec:standard}

The production acquisition is the mutual information between the response and the
latent firing rate \emph{at the query point itself},
%
\begin{equation}
  \Ustd(x^{*})
  \;=\; \Hmarg(x^{*}) - \E[\Hnoise(x^{*})]
  \;=\; H(R\mid x^{*},\mathcal D) - \E_{\rate\sim\text{post.}}\!\big[H(R\mid\rate)\big]
  \;=\; I\!\big(\rate(x^{*});R\mid x^{*},\mathcal D\big).
  \label{eq:Ustd}
\end{equation}
%
It scores ``how much would a spike count at $x^{*}$ pin down the firing rate at
$x^{*}$'' --- it depends on the marginal GP moments at $x^{*}$ only, with no $\px$
and no cross-covariance. It is exactly $\Uda$ with the predicted image taken to be
the query point, so the second term is the \emph{aleatoric} Poisson noise entropy
rather than an epistemic conditional entropy.

\subsubsection{The aleatoric noise entropy \texorpdfstring{$\E[\Hnoise]$}{E[H\_noise]} in closed form}
\label{subsubsec:hnoise}

Given a rate $\rate$, the entropy of $\Poi(\rate)$ is
$\Hnoise(\rate)=\rate\,(1-\log\rate)+\E_{r\sim\Poi(\rate)}[\log r!]$. Writing
$\lgr=\log\rate\sim\Gauss(\mug,\varg)$ and using the lognormal moment identities
%
\begin{equation}
  \E[e^{\lgr}] = \exp\!\big(\mug+\tfrac12\varg\big) = \fbar(x^{*}),
  \qquad
  \E[\lgr\,e^{\lgr}] = \exp\!\big(\mug+\tfrac12\varg\big)\,(\mug+\varg),
  \label{eq:lognormal-identities}
\end{equation}
%
(the first is the expected firing rate $\fbar$ of~\eqref{eq:fbar}), the first
piece integrates in closed form:
%
\begin{equation}
  \E_{\lgr}\!\big[\rate(1-\log\rate)\big]
  = \E_{\lgr}\!\big[e^{\lgr}(1-\lgr)\big]
  = \exp\!\big(\mug+\tfrac12\varg\big)\big(1-(\mug+\varg)\big)
  = -\exp\!\big(\mug+\tfrac12\varg\big)(\mug+\varg-1).
  \label{eq:hnoise-firstpiece}
\end{equation}
%
The remaining $\E[\log r!]$ is taken under the marginal predictive
$p(r\mid x^{*},\mathcal D)$ of~\eqref{eq:marg-predictive}. Hence
%
\begin{equation}
  \E[\Hnoise(x^{*})]
  \;=\; -\exp\!\big(\mug+\tfrac12\varg\big)(\mug+\varg-1)
        \;+\; \sum_{r} p(r\mid x^{*},\mathcal D)\,\log r! .
  \label{eq:Hnoise}
\end{equation}
%
This is the closed form used by \texttt{nd\_utility\_new} (its code comment cites
Eq.~33 of Goldin et al.\ 2023, PNAS). The production path
\texttt{standard\_utility} (i)~reads the marginal latent moments
$\pmean(x^{*}),\pvar(x^{*})$ via \texttt{get\_gp\_marginal\_moments},
(ii)~transforms them to $\mug,\varg$ by~\eqref{eq:logratemoments},
(iii)~optionally sets $\rmax$ adaptively, and (iv)~returns
$\Ustd=\Hmarg-\E[\Hnoise]$ from \texttt{nd\_utility\_new}.

% ---------------------------------------------------------------------------
\subsection{The distribution-aware (research) utility}
\label{subsec:da}

The distribution-aware utility~\eqref{eq:Uda} reduces uncertainty about firing
rates at the natural images we actually want to predict. It weights a candidate
$x^{*}$ by how informative its observation is about responses to typical natural
images --- i.e.\ by the cross-covariance between $x^{*}$ and images $x\sim\px$
(``distribution-aware'' $=$ aware of the natural-image distribution $\px$).
Evaluating $\Hcond$ needs the predictive moments of $\lat(x^{*})$ \emph{after a
hypothetical observation} $\lat(x)=\lat_i$ at a natural image $x$.

\subsubsection{Conditional GP moments via Gaussian conditioning}
\label{subsubsec:condmoments}

Let $\pv(x)=\Ktil^{-1}\kvec(x)$ be the SVGP projection and $q(\latt)=\Gauss(\vm,\vV)$
the variational posterior on the inducing values. The marginal latent posterior
moments at $x$ and $x^{*}$ are the projected moments of~\eqref{eq:predmoments},
$\pmean(x)=\pv(x)\transpose\vm$, $\pvar(x)=s(x)+\pv(x)\transpose\vV\pv(x)$
(and likewise at $x^{*}$), where $s(x)=\kvec(x,x)-\pv(x)\transpose\kvec(x)$ is
the prior Schur complement. The object that couples the two points is the
\textbf{prior conditional covariance} of $\lat(x),\lat(x^{*})$ given the inducing
values --- the term an earlier hand-assembled formulation omitted:
%
\begin{equation}
  \ccond \;=\; \kvec(x,x^{*}) - \pv(x)\transpose\Ktil\,\pv(x^{*}) .
  \label{eq:util-ccond}
\end{equation}
%
By the law of total covariance the full posterior cross-covariance is
%
\begin{equation}
  \Sigma_{x^{*}} \;=\; \ccond + \pv(x)\transpose\vV\,\pv(x^{*})
             \;=\; \kvec(x,x^{*}) + \pv(x)\transpose(\vV-\Ktil)\,\pv(x^{*}) ,
  \label{eq:util-crosscov}
\end{equation}
%
and standard Gaussian conditioning on the joint of
$(\lat(x),\lat(x^{*}))$ gives the conditional latent moments at $x^{*}$ after
observing $\lat(x)=\lat_i$,
%
\begin{align}
  \mu_{\mathrm{cond}}(x^{*})     &= \pmean(x^{*}) + \frac{\Sigma_{x^{*}}}{\pvar(x)}\,\big(\lat_i-\pmean(x)\big),
    \label{eq:util-mucond}\\[2pt]
  \sigma^{2}_{\mathrm{cond}}(x^{*}) &= \pvar(x^{*}) - \frac{\Sigma_{x^{*}}^{2}}{\pvar(x)} .
    \label{eq:util-sigcond}
\end{align}
%
These are transformed to log-rate space, $\gain\,\mu_{\mathrm{cond}}+\latz$ and
$\gain^{2}\sigma^{2}_{\mathrm{cond}}$, and fed to the same entropy
functional~\eqref{eq:Hmarg} to give the per-image conditional entropy
$H(R\mid x^{*},\lat(x),\mathcal D)$. The full Gaussian-conditioning derivation
(augmented inducing matrix, block inverse, and the equivalence to the rank-1
posterior update) is the deep-layer result~\eqref{eq:condmoments}; and the
conditional mean~\eqref{eq:util-mucond} obtained by conditioning on the joint
posterior coincides with the mean obtained by augmenting the inducing set with the
hypothetical observation and re-integrating (Thm.~\ref{thm:predmean-equiv}). Neither
is re-derived here.

\begin{remark}[The code is correct-by-construction]\label{rem:cond-correct}
\texttt{get\_gp\_conditional\_moments} builds the joint $[x;x^{*}]$, reads the
\emph{full} \texttt{covariance\_matrix} from the GPyTorch posterior, and takes
$\mu_{\mathrm{cond}},\sigma^{2}_{\mathrm{cond}}$ from its blocks
(with a clamp $\sigma^{2}_{\mathrm{cond}}\ge10^{-8}$). Because the variational
posterior covariance between any two points already equals
$\ccond+\pv\transpose\vV\pv^{*}=\Sigma_{x^{*}}$ --- the $\ccond$ term
included --- the code never forms the buggy $\pv\transpose\vV\pv^{*}$-only
cross-covariance. The ``missing-$\ccond$'' correction in the write-ups is a
warning about \emph{hand-assembled} cross-covariances (and about the deferred
\texttt{vargp\_direct} augmented-matrix path), not a defect in the shipped
predict path.
\end{remark}

\paragraph{Monte-Carlo estimator (one $\lat$ per image).}
Given pre-drawn samples $\{x_i\}_{i=1}^{N_{\mathrm{mc}}}\sim\px$:
%
\begin{enumerate}\itemsep2pt
  \item $\Hmarg(x^{*})=$ entropy~\eqref{eq:Hmarg} at the marginal moments
        $\mug,\varg$ of $x^{*}$.
  \item For $i=1,\dots,N_{\mathrm{mc}}$: read $(\pmean(x_i),\pvar(x_i))$; draw
        $\lat_i=\pmean(x_i)+\sqrt{\pvar(x_i)}\,\varepsilon,\ \varepsilon\sim\Gauss(0,1)$
        (if \texttt{sample\_lambda}; else $\lat_i=\pmean(x_i)$); form the
        conditional moments \eqref{eq:util-mucond}--\eqref{eq:util-sigcond} at
        $x^{*}$; set $\Hcond^{(i)}=$ entropy at $(\gain\mu_{\mathrm{cond}}+\latz,\ \gain^{2}\sigma^{2}_{\mathrm{cond}})$.
  \item $\Hcond=\tfrac{1}{N_{\mathrm{mc}}}\sum_i\Hcond^{(i)}$;\quad
        $\Uda=\Hmarg-\Hcond$.
\end{enumerate}
%
The $x_i$-posterior and the reparameterized draw $\lat_i$ are computed under
\texttt{torch.no\_grad()} (they do not depend on $x^{*}$, so excluding them avoids
$O(N_{\mathrm{mc}})$ graph accumulation); the conditional moments at $x^{*}$ are
differentiable in $x^{*}$. \texttt{sample\_lambda} is \emph{not} a debug flag:
$\texttt{False}$ uses the mean $\lat_i=\pmean(x_i)$ (a deterministic, reproducible
sanity check --- \S\ref{subsubsec:det}); $\texttt{True}$ is the correct unbiased
estimator (\S\ref{subsec:why-sample}).

\subsubsection{The deterministic special case}
\label{subsubsec:det}

Setting $\lat_i=\pmean(x)$ (the posterior mean) makes the innovation
$\lat_i-\pmean(x)=0$, so~\eqref{eq:util-mucond}--\eqref{eq:util-sigcond} collapse
to $\mu_{\mathrm{cond}}=\pmean(x^{*})$ (mean unchanged) and
%
\begin{equation}
  \sigma^{2}_{\mathrm{cond}}(x^{*}) \;=\; \pvar(x^{*})\,(1-\corr^{2}),
  \qquad
  \corr^{2} \;=\; \frac{\Sigma_{x^{*}}^{2}}{\pvar(x)\,\pvar(x^{*})},
  \label{eq:varreduce}
\end{equation}
%
where $\corr$ is the posterior latent correlation between $\lat(x)$ and
$\lat(x^{*})$. In log-rate coordinates the conditioning map is
$(\mug,\varg)\mapsto(\mug,\varg(1-\corr^{2}))$: same mean, variance scaled by
$1-\corr^{2}$. Writing $\Hmarg(\cdot,\cdot)$ for the spike-count
entropy~\eqref{eq:Hmarg} as a function of the log-rate moments, the utility is the
entropy difference
%
\begin{equation}
  \Uda(x^{*}) \;=\; \Hmarg(\mug,\varg) \;-\; \Hmarg\!\big(\mug,\ \varg(1-\corr^{2})\big).
  \label{eq:da-deterministic}
\end{equation}
%
So the deterministic $\Uda$ depends on exactly three quantities: the operating
point $\mug$, the marginal log-rate variance $\varg$ (how much entropy to start
with), and $\corr^{2}$ (the fraction of that variance conditioning removes).
$\corr^{2}$ is governed by the angle between $x^{*}$ and the conditioning image in
$\Cmat$-space: a small angle gives $\corr^{2}\!\to\!1$ (near-total variance
removal), a large angle $\corr^{2}\!\to\!0$ (no conditioning). This same $\corr^{2}$
sets the sampling bias of \S\ref{subsec:why-sample}.

% ---------------------------------------------------------------------------
\subsection{Why we sample the latent instead of using its mean}
\label{subsec:why-sample}

$\Hcond$ is a double expectation (Remark~\ref{rem:double-exp}), over $x\sim\px$ and
over $\lat(x)\sim q$. The tempting shortcut replaces the inner expectation by
evaluation at the mean,
%
\begin{equation}
  \widetilde{H}_{\mathrm{cond}}(x^{*})
  \;=\; \E_{x\sim\px}\!\big[\,H\big(R\mid x^{*},\lat(x)=\pmean(x),\mathcal D\big)\big].
  \label{eq:mean-shortcut}
\end{equation}

\begin{proposition}[The mean shortcut is biased]\label{prop:bias}
$\E_{\lat(x)}[H(R\mid x^{*},\lat(x),\mathcal D)]\neq H(R\mid x^{*},\lat(x)=\pmean(x),\mathcal D)$
in general. The predictive mean of $\lat(x^{*})$ is linear in the conditioning
value, $\E[\lat(x^{*})\mid\lat(x),\mathcal D]=\pmean(x^{*})+\alpha\,(\lat(x)-\pmean(x))$
with regression coefficient $\alpha=\mathrm{Cov}(\lat(x^{*}),\lat(x)\mid\mathcal D)/\pvar(x)$,
while the predictive \emph{variance} does not depend on the value of $\lat(x)$.
Entropy is nonlinear in the predictive mean, so a second-order Taylor expansion of
$H$ about $\lat(x)=\pmean(x)$, with $\E[\lat(x)-\pmean(x)]=0$, leaves the leading
bias
%
\begin{equation}
  \E_{\lat(x)}[H] - H\big|_{\lat(x)=\pmean(x)}
  \;\approx\; \tfrac12\,\alpha^{2}\,\pvar(x)\;\frac{\partial^{2}\Hmarg}{\partial\pmean^{2}}\Big|_{x^{*}} .
  \label{eq:bias-raw}
\end{equation}
%
Since $\mathrm{Cov}(\lat(x^{*}),\lat(x))=\corr\,\sqrt{\pvar(x^{*})\pvar(x)}$, we
have $\alpha^{2}\,\pvar(x)=\corr^{2}\,\pvar(x^{*})$, and the $\pvar(x)$ cancels:
%
\begin{equation}
  \mathrm{Bias} \;\approx\; \frac{\corr^{2}}{2}\;\pvar(x^{*})\;\frac{\partial^{2}\Hmarg}{\partial\pmean^{2}}\Big|_{x^{*}} .
  \label{eq:bias}
\end{equation}
\end{proposition}

The bias vanishes \emph{only} when $\corr=0$ ($x$ uninformative about $x^{*}$ ---
but then the utility contribution is $\approx0$ anyway), or $\pvar(x^{*})=0$
(nothing left to learn), or $\partial^{2}\Hmarg/\partial\pmean^{2}=0$ (entropy
linear in the mean --- false for the Poisson-lognormal). Notably $\pvar(x)=0$ is
\emph{not} a zero-bias condition: it cancels in~\eqref{eq:bias}.

\begin{remark}[The bias is largest exactly where it matters]\label{rem:bias-matters}
Bias $\propto\corr^{2}\,\pvar(x^{*})$ is maximal when the sampled image $x$ is most
correlated with $x^{*}$ (most relevant to the utility) \emph{and} when $x^{*}$ is
most uncertain (where the acquisition should be steering us) --- i.e.\ largest as
$\corr\to1$. A systematic bias does not shrink with more Monte-Carlo samples: it
converges to the \emph{wrong} value and corrupts the optimizer's direction.
Sampling $\lat$ is unbiased and its variance falls as $O(1/N)$. Hence: sample
$\lat$.
\end{remark}

\paragraph{One $\lat$ per image is optimal (nested MC).}
With $g(x,\lat)=H(R\mid x^{*},\lat,\mathcal D)$, let $\tau^{2}=\mathrm{Var}_x[\E_{\lat}g]$
(between-image) and $\E_x[\sigma^{2}_x]$ (within-image). For a fixed budget of $M$
entropy evaluations, using $N_{\lat}$ latent samples per image gives variance
$(N_{\lat}\tau^{2}+\E_x[\sigma^{2}_x])/M$, strictly worse for $N_{\lat}>1$.
Spending the budget on \emph{more images}, one $\lat$ each, is optimal --- which is
what the implementation does.

\caveat{The ``practical magnitude'' analysis in the source
(\texttt{why\_sample\_lambda.tex}) is self-flagged ``review carefully.'' Pushed to
log-rate space the curvature carries a gain factor,
$\partial^{2}\Hmarg/\partial\pmean^{2}=\gain^{2}\,\partial^{2}\Hmarg/\partial\lgr^{2}$,
so $\mathrm{Bias}\propto\gain^{2}$; for a small gain (this codebase typically has
$\gain\approx0.03$) the bias is far below the MC standard error, and the
deterministic and sampled $\Uda$ nearly coincide. The regime guide $\gain\ll1$
negligible / $\gain\sim1$ moderate / $\gain\gg1$ essential, and the specific
numbers, are plausibility arguments, not settled results. The core claim ---
bias~\eqref{eq:bias} does not vanish with $N$ --- is solid; the magnitude estimate
is provisional. The \emph{correct} estimator remains the sampled one regardless.}

% ---------------------------------------------------------------------------
\subsection{The entropy landscape over image space}
\label{subsec:landscape}

\paragraph{Monotone shape --- utility is not pure epistemic uncertainty.}
The spike-count entropy $\Hmarg(\mug,\varg)$ increases monotonically in
\emph{both} arguments: with $\mug$ (more predicted firing) and with $\varg$ (more
GP/epistemic uncertainty). A candidate therefore scores high for a \emph{mixture}
of reasons --- it is predicted to fire a lot \emph{and} it is epistemically
uncertain. The utility is thus not a pure epistemic-uncertainty acquisition;
predicted firing rate is baked in. This is a deliberate modelling choice: a
normalized arc-cosine kernel would force $\Kf(x,x)=1$, killing the norm channel so
that utility tracked only angular/RF alignment (pure epistemic), but held-out
$\text{test\_r}$ then drops by $\sim$25\% (a fitted PNAS cell, $\Mind=100$: $0.79\!\to\!0.59$)
--- image norm is genuinely informative for neural encoding, so the production
kernel keeps the norm dependence.

\paragraph{The $\rmax$ truncation artifact.}
With a fixed $\rmax=100$ the Laplace sum~\eqref{eq:Hmarg} misses the probability
mass above $\rmax$ once the Poisson peak climbs past it, and $\Hmarg$
\emph{collapses to $\approx0$} for $\mug\gtrsim\log100\approx4.6$ (and at lower
$\mug$ for large $\varg$). This is a numerical artifact, not real entropy: the
computable region is roughly triangular in $(\mug,\varg)$. A cheap pre-check
without running the Laplace sum is $z_{\mathrm{safe}}=(\log\rmax-\mug)/\sqrt{\varg}$
--- $z_{\mathrm{safe}}>2$ safe, $<1.5$ unreliable, $<0.5$ broken. Natural images
from the pool sit at $z_{\mathrm{safe}}\gg10$, always safe; trouble arises only for
artificially amplified images. The \textbf{adaptive-$\rmax$ guard}
(\S\ref{subsec:code}) removes the collapse.

\paragraph{Norm-exploitation divergence (research path).}
Because the arc-cosine kernel is positively $1$-homogeneous
($\Kf(\scal x,y)=\scal\,\Kf(x,y)$ when the bias $\sigz=0$), scaling
$x^{*}\to\scal x^{*}$ scales $\pmean(x^{*})\sim\scal$ and $\pvar(x^{*})\sim\scal^{2}$,
so $\Hmarg$ grows while a directionally aligned conditioning image keeps
$\corr\approx1$ and $\sigma^{2}_{\mathrm{cond}}\approx0$. Under free-space gradient
ascent $\Uda$ then grows $\sim\scal^{1.9}$ (numerically) --- the optimizer
``diverges'' by amplifying image norm rather than finding angularly informative
images. This is intrinsic to the homogeneous kernel, \emph{not} a bug; the full
statement and the $\sigz^{2}>0$ damping (alignment peak strictly at $\scal=1$) are
Thm.~\ref{thm:divergence} in the deep layer. Over the finite natural-image
\emph{pool} (production selection over discrete candidates) norms are bounded and
this never fires.

\caveat{\textbf{Utility accuracy ceiling (known limitation).} At high firing the
two utilities fail in opposite directions through the same $\rmax=100$ truncation:
$\Ustd$ can run to $+\infty$ while $\Uda$ silently collapses toward $0$. The firing
-rate guard ($f_{\max}$, default $100$) keeps generation in the accurate regime; a
dedicated DA-path look-up table was deliberately not built. Present as a limitation
of the fixed-$\rmax$ path, cured on the standard path by the adaptive guard.}

% ---------------------------------------------------------------------------
\subsection{Code realization and conventions}
\label{subsec:code}

\begin{table}[t]
  \centering
  \footnotesize
  \begin{tabular}{@{}lll@{}}
    \toprule
    Function & Input moment convention & Sum index\\
    \midrule
    \texttt{compute\_H}$(\mu,\sigma^{2},\rmax,\gain,\latz)$ & \emph{raw} $\lat$ moments (transformed inside) & $r=0,\dots,\rmax-1$\\
    \texttt{nd\_utility\_new}$(\mug,\varg,\rmax)$ & \emph{log-rate} $\lgr$ moments (pre-transformed) & $r=0,\dots,\rmax$\\
    \bottomrule
  \end{tabular}
  \caption{The two entropy entry points differ in \emph{both} the moment
  convention and the truncation range. $\Ustd$ transforms manually then calls
  \texttt{nd\_utility\_new}; $\Uda$ passes raw moments $+(\gain,\latz)$ to
  \texttt{compute\_H}.}
  \label{tab:code-conventions}
\end{table}

\paragraph{Entry points.} $\Ustd$ runs \texttt{standard\_utility}, which calls
\texttt{nd\_utility\_new}. The research routine
\texttt{distribution\_\allowbreak aware\_\allowbreak utility} computes $\Hmarg$
with \texttt{compute\_H}, and per MC sample forms the conditional moments
(\S\ref{subsubsec:condmoments}) and calls \texttt{compute\_H} again for
$\Hcond^{(i)}$. Its marginal and conditional latent moments come from the two
\texttt{get\_gp} moment helpers of \S\ref{subsubsec:condmoments}.

\paragraph{The adaptive-$\rmax$ guard.} \texttt{compute\_adaptive\_rmax} sizes the
truncation from the moments: with an upper log-rate
$\log\rate_{\uparrow}=\mug+k\sqrt{\varg}$ ($k=3$ by default),
$\rmax=\lceil\rate_{\uparrow}+5\sqrt{\rate_{\uparrow}}+10\rceil$ (a $5$-SD Poisson
margin), floored at $\text{min\_rmax}$ (default $200$). It \emph{raises} rather
than silently truncating when the needed $\rmax$ would exceed
$\text{max\_rmax}$ (default $10000$) --- an aggressive-truncation guard, since a
wrong entropy would otherwise pass unnoticed. This is what keeps $\Hmarg$ from
collapsing when $\mug$ is high; non-adaptive callers pass a fixed $\rmax$
(production default $100$).

\paragraph{Gradient compatibility.} Every production entropy/utility function is
differentiable in $x^{*}$: the Laplace path uses \texttt{torch.where} (not indexed
assignment, which breaks autograd), and the Lambert~$W$ has a proper custom
backward $dW/dy=W/(1+W)$. The caller chooses to track gradients (pass
\texttt{requires\_grad} tensors, enabling gradient-ascent $x^{*}$ optimization ---
and exposing the divergence above) or to suppress them (\texttt{torch.no\_grad()}).

\caveat{\textbf{A/$\latz$ convention split.} \texttt{compute\_H} takes \emph{raw}
$\lat$ moments and applies $\mug=\gain\pmean+\latz,\ \varg=\gain^{2}\pvar$
internally, whereas \texttt{nd\_utility\_new} and the low-level
\texttt{\_diff\_laplace\_log\_probs} take \emph{already-transformed} $\lgr$ moments
(Table~\ref{tab:code-conventions}). Both deliver the same math to the Laplace
pipeline, but silently passing raw moments where $\lgr$ moments are expected (or
vice versa) is a genuine bug source --- keep the mapping explicit at every call.}

\caveat{\textbf{$\rmax$ off-by-one.} \texttt{compute\_H} sums $r=0,\dots,\rmax-1$
while \texttt{nd\_utility\_new} sums $r=0,\dots,\rmax$. Both are finite truncations
of~\eqref{eq:Hmarg}; the one-bin difference is immaterial at low firing but is a
real inconsistency between the two entry points --- do not treat $\rmax$ as
bit-identical across them.}

\caveat{\textbf{Deprecated helper.} The current local \texttt{utils.py} has
\emph{no} \texttt{lambda\_moments}: that name is a legacy varGP routine, living in
the deprecated repo-root \texttt{utils.py}. Its role is now filled by
\texttt{get\_gp\_marginal\_moments} (the marginal moments $\pmean,\pvar$) and
\texttt{get\_gp\_conditional\_moments} (the conditional pair
$\mu_{\mathrm{cond}},\sigma^{2}_{\mathrm{cond}}$). Do not resurrect it.}

\begin{remark}[Mode restriction and an analysis-only routine]\label{rem:mode-restriction}
$\Uda$ currently requires a model exposing \texttt{model(X).\allowbreak covariance\_matrix}
(the \texttt{default\_gpy} GPyTorch VariationalGP); the eigenspace
\texttt{vargp\_direct} path is deferred (it would need the augmented-matrix conditioning
of~\eqref{eq:condmoments}). Separately, \texttt{compute\_H\_MC} (sample
$\lgr\sim\Gauss$, $r\sim\Poi(e^{\lgr})$, average the clipped Laplace $\log p(r)$)
exists \emph{only} for pedagogy/analysis: it is non-differentiable (discrete
Poisson sampling) and biased (clipping), and carries an in-code
``\textsc{raise to user}'' warning about hard-coded values --- it is never a
production acquisition.
\end{remark}
